\documentclass[11pt]{article}

\usepackage[T1]{fontenc}
\usepackage[utf8]{inputenc}
\usepackage{lmodern}
\usepackage[margin=1in]{geometry}
\usepackage{amsmath}
\usepackage{amssymb}
\usepackage{booktabs}
\usepackage{hyperref}
\usepackage{enumitem}

\newcommand{\todoitem}[1]{%
  \par\noindent\textbf{TODO:} #1\par
}

\title{Reasonable Quaternion Arithmetic for \texttt{polyForth} on Volatco\\
\large Representation, Implementation, and Performance Evidence}
\author{Cartheur}
\date{\today}

\begin{document}
\maketitle

\begin{abstract}
This paper studies whether quaternion arithmetic is a reasonable computational
choice for \texttt{polyForth} on Volatco, an asynchronous F18A-based platform
intended for low-power control, robotics, and experimental machine-intelligence
workloads. The standard of reasonableness adopted here is deliberately narrow:
the representation must be simple to encode in a stack language, compact enough
to respect the constraints of an 18-bit node, and measurably competitive with
plausible alternatives on representative rotational workloads. The repository
therefore combines source code, benchmark definitions, and manuscript material
so that the argument can be revised against hardware measurements rather than
asserted from algebra alone.
\todoitem{Revise the abstract once real Volatco measurements exist so that it
states actual findings rather than the present research intent.}
\end{abstract}

\section{Introduction}

Quaternion methods are attractive in rotational computation because they offer a
compact composition law and avoid some of the representational overhead of
matrix-based formulations. Those advantages are familiar in conventional
software environments, but they cannot simply be assumed on a tiny asynchronous
multicomputer running \texttt{polyForth}. On such a platform, the right
question is not whether quaternions are elegant, but whether they are
operationally reasonable.

Volatco is a useful target for this question because it is explicitly
Forth-native and designed around deterministic, low-latency execution. Its
published documentation emphasizes reproducible \texttt{polyForth} experiments
with explicit pass/fail criteria, power measurement points, and bounded-latency
behavior \cite{volatco-site,volatco-docs}. That combination makes it possible
to write a paper that is both mathematical and empirical.

\todoitem{Add citations from the local quaternion corpus for the mathematical
motivation, especially references on quaternion rotation and computational
advantages or limitations relative to matrices.}

\section{Claim}

The working claim of this paper is modest:

\begin{quote}
A fixed-point quaternion representation is a reasonable rotational algebra for
\texttt{polyForth} on Volatco when it yields compact code, controlled state,
and reproducible performance under explicitly measured workloads.
\end{quote}

The claim is intentionally narrower than ``quaternions are best.'' It is a
claim about suitability under a specific architecture, runtime style, and
evidence standard.

\todoitem{Replace the current claim text with a version tied to measured
thresholds, for example code size, elapsed time, or power bounds, once those
numbers exist.}

\section{Representations Under Comparison}

The implementation scaffold in this repository uses two competing
representations. The quaternion kernel stores $(a,b,c,d)$, where $a$ is the
scalar part and $b,c,d$ are the vector components. The comparison kernel stores
a $3\times 3$ rotation matrix in row-major order plus a three-component vector.
This side-by-side structure is important: the paper should not ask whether
quaternions are elegant in the abstract, but whether they remain reasonable
once compared to a direct matrix approach under the same tooling and target
constraints.

The quaternion layout maps directly to Hamilton's multiplication rule. The
current draft keeps the representation abstract enough to allow either integer
arithmetic or scaled fixed-point arithmetic after target measurements determine
which is more practical on Volatco.

\todoitem{Choose the actual numeric model for the paper: plain integer,
fixed-point with explicit scale factor, or a mixed representation. Document the
overflow and precision tradeoffs for the chosen model.}

The scalar-first ordering is not mathematically necessary. It is adopted here
for engineering clarity, because conjugation, norm-squared, and scalar-vector
decomposition can be written with simple address-based stack effects in Forth.

\todoitem{Add one short worked example showing the in-memory layout of a
quaternion and a matrix on the target Forth system.}

\section{Implementation Strategy}

The initial Forth implementation favors transparency over optimization. Core
quaternion words are provided for addition, subtraction, conjugation,
norm-squared, Hamilton product, and vector rotation via
$qvq^{-1}$, implemented as \texttt{q * v * conjugate(q)}. A companion matrix
kernel performs matrix-vector and matrix-matrix multiplication in the same
address-based style. This is the correct place to start because a paper cannot
defend an optimized kernel unless it first presents readable reference
implementations for both the proposed and comparison methods.

\todoitem{Add an appendix excerpt of the actual reference Forth code and note
any edits required for the exact Volatco \texttt{polyForth} image.}

After reference behavior is established, optimization can proceed in three
directions:

\begin{enumerate}[label=\arabic*.]
\item reduce stack traffic
\item remove scratch storage where the target system permits
\item distribute work across nodes when the communication cost is justified
\end{enumerate}

\todoitem{State whether this paper evaluates only single-node kernels or also
includes multi-node decomposition across the Volatco mesh.}

\section{Measurement Plan}

The benchmark argument in this repository compares quaternion operations against
a direct $3\times 3$ rotation matrix kernel written in the same style. That
comparison is necessary because the matrix method spends more storage but may
win in some workloads by avoiding quaternion conjugation or normalization
overhead.

Each benchmark should record:

\begin{itemize}
\item iteration count
\item elapsed time
\item code footprint
\item current or power measurement
\item watchdog stability during long runs
\end{itemize}

The results should be treated as evidence about suitability for Volatco, not as
universal claims about all Forth systems.

\todoitem{Insert a real results table here with at least quaternion rotation
versus matrix-vector rotation, including iteration count, elapsed time, and
power measurement point.}

\todoitem{Document the timing method precisely: serial timestamping, on-board
counter, external logic capture, or another method.}

\section{Expected Contribution}

If the measurements support the working claim, the paper contributes a concrete
case study in how a classical algebraic formalism can be adapted to a modern
asynchronous Forth-native platform. If the measurements do not support the
claim, that result is still valuable because it clarifies where quaternion
methods cease to be worth their conceptual appeal under tight embedded
constraints.

The comparison with a matrix kernel is especially important here. A paper that
only demonstrates that quaternion code runs on Volatco would be incomplete. The
interesting result is whether quaternion state, composition, and rotation
support a better tradeoff than a simpler matrix pipeline once timing, code
size, and power are measured on the same board.

\todoitem{Once data exists, make this section argue from the observed tradeoff
rather than from expected contribution alone.}

\section{Current Status}

At present, the repository contains:

\begin{itemize}
\item an address-based quaternion kernel for Forth
\item a companion $3\times 3$ matrix kernel
\item a benchmark scaffold for repeated arithmetic and rotation workloads
\item a reproducibility protocol for eventual Volatco measurements
\end{itemize}

It does not yet contain:

\begin{itemize}
\item confirmed target-specific \texttt{polyForth} adaptations
\item hardware measurements from a real Volatco board
\item a completed citation apparatus grounded in the local quaternion corpus
\end{itemize}

The present engineering assumptions about Volatco derive from the public site
and documentation rather than from private hardware notes
\cite{volatco-site,volatco-docs}.

\todoitem{Replace this status section with a brief methodology section after the
first hardware campaign is complete.}

\section{Next Revision Points}

\begin{itemize}
\item replace placeholder prose with citations from the local quaternion corpus
\item confirm \texttt{polyForth} compatibility on Volatco
\item choose and document a fixed-point scaling rule
\item add measured benchmark results
\item interpret the quaternion results against the matrix baseline
\item narrow the claim to match the data
\end{itemize}

\section{Focused TODO List}

\begin{itemize}
\item Gather mathematical citations from the local quaternion papers and books.
\item Confirm the exact Volatco \texttt{polyForth} dialect and any source edits
needed for the current Forth files.
\item Choose the numeric representation and document scale, range, and overflow
behavior.
\item Run quaternion and matrix benchmarks on real Volatco hardware.
\item Add tables for timing, power, and code size.
\item Rewrite the abstract, claim, and discussion around measured evidence.
\end{itemize}

\bibliographystyle{plain}
\bibliography{references}

\appendix
\section{Benchmark Protocol Summary}

The benchmark protocol should report, at minimum, board revision, runtime
revision, exact source, iteration count, elapsed time, power measurement point,
code size, pass/fail outcome, and notes on watchdog or reset behavior. The
baseline workload set should include repeated quaternion addition, Hamilton
product, norm-squared, vector rotation, matrix-vector rotation, and
matrix-matrix multiplication.

\end{document}
